% Part: first-order-logic
% Chapter: introduction
% Section: soundness-completeness

\documentclass[../../../include/open-logic-section]{subfiles}

\begin{document}

\olfileid{fol}{int}{scp}

\olsection{السلامة والاكتمال}

ومن مباحثنا أيضًا أنساق ال!!{derivation} لمنطق الرتبة الأولى.
وللمناطقة فيها أنساق !!{derivation} كثيرة، تتفق كلها في علاقة
ال!!{derivability} التي تقيمها بين ال!!{sentence}s. نقول بإمكان
\emph{!!{derive}}~$!A$ من ${\OLId{greek-variable}{Gamma}}$، ونرمز
${\OLId{greek-variable}{Gamma}} \Proves !A$، متى وجد !!a{derivation} من النوع المحدد
بالتعريف. وال!!^{derivation}s ترتيبات منتهية من الرموز، كقائمة من
ال!!{sentence}s أو تركيب أعقد من القائمة. وإنما نطلب من أنساق
ال!!{derivation} وسيلة للفصل في كون !!a{sentence} لازمة عن
مجموعة~${\OLId{greek-variable}{Gamma}}$. فلا تفي بهذا الغرض إلا إذا كان
${\OLId{greek-variable}{Gamma}} \Entails !A$ مكافئًا لـ${\OLId{greek-variable}{Gamma}} \Proves !A$.

فإن تحقق ${\OLId{greek-variable}{Gamma}} \Proves !A$ دون ${\OLId{greek-variable}{Gamma}} \Entails !A$،
تجاوز نسق ال!!{derivation} الحد، وأثبت ما لا يلزم.
وكون ${\OLId{greek-variable}{Gamma}} \Entails !A$ لازمًا كلما كان
${\OLId{greek-variable}{Gamma}} \Proves !A$ هو \emph{السلامة}، وهي أقل ما نطلب من
نسق !!{derivation} صالح. وعلى العكس، إن تحقق
${\OLId{greek-variable}{Gamma}} \Entails !A$ دون ${\OLId{greek-variable}{Gamma}} \Proves !A$، قصر نسق
ال!!{derivation} عن إثبات بعض ما يلزم. وكون
${\OLId{greek-variable}{Gamma}} \Proves !A$ متحققًا كلما كان ${\OLId{greek-variable}{Gamma}} \Entails !A$
هو \emph{الاكتمال}. وبرهان السلامة أيسر عادة؛ نستقرئ فيه تركيب
ال!!{derivation}s، إذ هي معرفة استقراءً. وأما الاكتمال فبرهانه
أشد مؤونة.

وتتفرع عن السلامة والاكتمال نتائج مهمة. إذا أمكن !!{derive}
تناقض، مثل $!A \land \lnot !A$، من مجموعة ال!!{sentence}s~${\OLId{greek-variable}{Gamma}}$،
وصفت المجموعة بأنها \emph{غير متسقة}. ولا نموذج لمجموعة
غير متسقة~${\OLId{greek-variable}{Gamma}}$، فهي لا تقبل الإرضاء. ويعطي الاكتمال
العكس: كل ${\OLId{greek-variable}{Gamma}}$ ليست غير متسقة، أي هي \emph{متسقة}،
لها نموذج. وهذه العبارة تكافئ الاكتمال، وهي التي سنبرهنها
في الحقيقة. وفيها نتيجة عميقة قد تستوقف القارئ: عدم إمكان إثبات
$!A \land \lnot !A$ من ${\OLId{greek-variable}{Gamma}}$ يكفل وجود
!!a{structure} على ما تصفه~${\OLId{greek-variable}{Gamma}}$.
فالجواب عن سؤال «أي مجموعات ال!!{sentence}s لها نماذج؟»
هو: المتسقة كلها، ولا شيء سواها.

ومن أهم النتائج مبرهنة التراص ومبرهنة لوفنهايم--سكولم. لهما
مكان عظيم في نظرية النماذج، وتستخرج منهما نتائج كثيرة؛ وقد
تقدم مثالان: امتناع التعبير بمنطق الرتبة الأولى عن تناهي
!!{domain} !!a{structure}، وعن كون مجموعة المجال
!!{nonenumerable}.

لم يحصل هذا الضبط دفعة واحدة في تاريخ المنطق؛ بل طال النظر
في تعريف النحو والدلالة، ووضع أنساق !!{derivation} الصالحة،
وبرهنة سلامتها واكتمالها، وتبيين ما تعبر عنه لغات الرتبة الأولى
وما لا تعبر عنه. وقد عرفنا اليوم كيف نفي بهذه المطالب،
لكن تتبع التفاصيل لا يخلو من مشقة أو التباس. وفائدته باقية:
فالطرائق التي نبنيها هنا للغة منطق الرتبة الأولى تستعمل في
سائر المنطق، وفي علوم الحاسوب واللسانيات. فالعناية بهذه
التفاصيل عون على ما بعدها.

\end{document}
